\subsection{Grundsätze für menschenzentriertes Design}
\label{ch:SdF:UI:sec:menschenzentriertes_Design}

Bei der Entwicklung von neuen Systemen, wird meist erst die Technologie designt. Vor allem in diesen Technologie-fokussierten Systemen werden durch deren Komplexität und Unintuitivität Nutzungsfehler gemacht. Wird das Design \textit{menschenzentriert} erarbeitet, so werden zuerst die Anforderungen an die Maschine erfasst \cite{Vollrath:2015:Taschenbuch, Wickens:1995}.

% *******************************************************************
%           Fokus auf Intuition
% *******************************************************************
Der Fokus liegt auf der \textit{Intuition} des Systems. Diese beschreibt, dass ein Nutzer den Umgang mit einem System schnell erlernen kann und im Zusammenhang damit ohne eine umfangreiche Erklärung versteht, wie er dieses zu verwenden oder zu bedienen hat. Ein Beispiel ist die Zusammenfassung komplexer Systemdaten und -werte auf eine automatisierte Einschätzung sowie Ausgabe zu dessen Zustand \cite{Vollrath:2015:Taschenbuch, Wickens:2013}.

% *******************************************************************
%           visuelle Suche + Pop-Out-Effekt
% *******************************************************************
Die Darstellung von Systemdaten sollte so erfolgen, dass sie gut lesbar sind und intuitiv erkennbar ist, was dargestellt ist.
% Area of Interest + visuelle Aufmerksamkeit
Um sich auf bestimmte Bereiche zu konzentrieren, muss die eigene \textit{Aufmerksamkeit} darauf gelenkt werden. Dabei ist die \textit{Area of Interest} (dt. der Bereich des Interesses) nur so groß, wie diese mit den Augenbewegungen abgedeckt werden kann \cite{Vollrath:2015:Taschenbuch, Wickens:2013}.

% ****************
% visuelle serielle suche + extensive = fatiguing
Werden auf einer Informationsanzeige viele Werte dargestellt, so muss nach dem gewünschten Ergebnis \textit{seriell} gesucht werden. Diese monotone Aufgabe bedeutet, dass alle Ergenisse stückweise gescannt werden.
Nutzer die für die Überwachung der Daten zuständig sind, müssen dennoch jederzeit auf ein unwahrscheinliches und entsprechend unerwartetes Event reagieren können. Dabei kann diese anstrengende und umfassende Suche, auch \textit{extensive scanning} genannt, nur für eine gewisse Zeit erfolgen.
% parallere Suche + Pop-Out (Farben, Formen, etc.)
Im Unterschied dazu steht die \textit{parallele Suche}, welche eine Art scannen auf einen Blick ermöglicht. Um dies zu ermöglichen, müssen gesuchte Ziele eindeutig von ihrer Umgebung unterscheidbar sein. Während Objekte mit beispielweise der selben monochromen Farbe dazu tendieren, ähnlich verarbeitet zu werden, stechen Objekte durch diesen \textit{Pop-Out-Effekt} hervor. Dieser kann auf Grundlage von Form, Farbgebung, -intensität, -sättigung oder Kontrast entstehen. Die Unterscheidung durch Farbsättigung zeigt dabei beispielsweise eine größere Wirkung als unterschiedliche Farben gleicher Sättigung. Bis auf Farben mit etablierter Bedeutung, wie rot als Signalfarbe, ist es kaum möglich, die Absicht zu deuten.
% Suchzeiten und Ablenkungen -> heterogener clutter
Die Suchzeiten nach einem Ziel erhöhen sich deutlich, wenn die Anzahl der Ablenkungen steigt, bleiben jedoch gleich, wenn die Unterscheidbarkeit zwischen Ziel und Ablenkung groß ist.
Während ein homogenes Hintergrundgeschehen und eine einzigartige Eigenschaft des Ziels die Suche erleichtern, verlängert sich die Suchzeit, wenn die Anzeige mit Daten überfüllt ist und die Umgebung mit heterogenem clutter (dt. Gerümpel) gefüllt ist.
Es muss ein Kompromiss zwischen der Suchzeit und Genauigkeit eingegangen werden. Genauere Suchen scheinen langsamer zu sein, wobei allerdings Faktoren wie die Gleichheit zwischen Ziel und Ablenkung ebenso dazu neigen, Fehler zu verursachen \cite{Wickens:2013, Vollrath:2015:Taschenbuch}.

% ****************
% Choice of Action - TODO: eher egal, hab eh zu viel Text
%...

%--> visuelle Aufmerksamkeit\\
%    ----> selektive Aufmerksamkeit + Entscheidungen\\
%    ----> good decision making\\
%\cite{Wickens:2013}
%- welche Parts zu verarbeiten, welche auszusortieren
%- viele Entscheidungen sind iterativ
%-> initiale Hypothese / Einstellung wird die Suche nach weiteren Informationen triggern
%-> choice of action
%(-> wenn diese choice eindeutig durch farbgestaltung abgenommen wird)
%-	Good decision making
%-> Werte sind oft persönlicher Natur -> macht es zur schwierigen Aufgabe
%-> selbst wenn eine Einigung besteht, ist die optimale Entscheidung (bei Wiederholung; mit genügend Zeit dafür) oft nicht optimal für die einzelne Wahl
%--> vor allem unter Zeitdruck - ohne Zeit zur vollständigen Diagnose und Berücksichtigung aller Ergebnisse (Wickens et al. 2013, S. 248–249)

% ****************
% Prinzipien
Wickens erwähnt fünf Prinzipien, welche dabei unterstützen, nutzerfreundliche Dashboards zu entwerfen. Die Darstellung hat den Anforderungen der Aufgabe zu entsprechen. Darüber hinaus sollen dem Nutzer mentale Aufgaben, wie das Zusammenzählen von Werten für ein Gesamtbild, abgenommen werden. Die Graphen sollen eindeutig interpretiert werden können, wobei die \textit{Data-Ink-Ratio} (dt. Verhältnis von veranschaulichten Daten zu angezeigter Farbe) hoch zu sein hat. Für den Überblick sollen unterschiedliche Graphen weiterhin konsistent aufgebaut werden.
% Konsistenz + Design Trade-Offs
Werden mehrere Graphen dargestellt, so ist es wichtig, deren Unterschiede in Form einer Legende oder der Symbolik selbst herauszuarbeiten. Um jedoch dieses konsistente Bild über das gesamte System zu erzeugen, muss ein Kompromiss gefunden werden, sodass nicht immer alle der fünf Prinzipien angewendet werden können. Damit ist beispielsweise die konsistente Nutzung von Farben, Symbolen, Formen oder Bedienelementen gemeint \cite{Vollrath:2015:Taschenbuch, Wickens:2013}.

% *******************************************************************
%           Gesetz der Nähe
% *******************************************************************
Werden verschiedene Graphen häufig in einem gemeinsamen Kontext verwendet, so besagt das \textit{proximity compatibility principle} (dt. Gesetz der Nähe), dass diese zusammengehörigen Informationen beisammen darzustellen sind.
Bei der Suche nach Informationen von verschiedenen Oberflächen eines Systems muss zusätzlich zu der Suchzeit auf einer Oberfläche die Zeit zum Wechseln zwischen und Finden der gewünschten Oberfläche einberechnet werden
\cite{Wickens:1995, Wickens:2013}.

% *******************************************************************
%           Konzept: Lernen + Fehler
% *******************************************************************
Menschen lernen anhand von Wiederholungen. Je öfter sie eine Aufgabe ausführen oder ein System bedienen, umso leichter fällt es ihnen. Je intuitiver ein System ist, desto weniger Lernaufwand wird benötigt. Gleichzeitig fällt es jüngeren Menschen leichter, sich in einer Multi-Tasking-Aufgabe auf einen bestimmten Aspekt zu konzentrieren. Dafür wird die Leistung des Kurzzeitgedächtnisses sowie ein gewisser Übungsgrad benötigt. Ältere Menschen sind nicht so effizient und es fällt ihnen schwerer, diesen Aspekt in den vielen gleichzeitigen Aufgaben zu fokussieren. Somit wird das zuvor beschriebene \textit{extensive scanning} noch anspruchsvoller.
% Fehler
\glqq Irren ist menschlich\grqq{}\footnote{\url{https://de.wiktionary.org/wiki/Irren_ist_menschlich} [Zugegriffen am: 25.08.2025]} hat Seneca gesagt. Während Fehler unterschiedlichster Art beabsichtigt oder nicht passieren, können \textit{Slips} (dt. Ausrutscher oder Überseher) in der Überwachung von Datenlagen problematisch werden. Laut Defintion sind es keine Fehler, da die Intention und das Verhalten des Nutzers korrekt sind, jedoch führen sie zur Ausführung falscher Aktionen. Wenn eine Aufgabe automatisiert abläuft, liegt kaum Aufmerksamkeit auf der Ausführung, wodurch der Nutzer seinen Fehler nicht einmal mitbekommt \cite{Wickens:2013}.

% *******************************************************************
%           Automation: Anti-Lernen
% *******************************************************************
Zur Vermeidung von Fehlern werden Aufgaben automatisiert. Für Aufgaben, die Menschen nicht, nur mit schlechter Qualität oder hohem Aufwand lösen können, bieten Maschinen effiziente Lösungen.
Dabei soll der Mensch mit allen nötigen Informationen versorgt werden. Um eine Automation menschenzentriert zu gestalten, sollte dessen Ergebnisse verständlich und Vorgänge leicht zugänglich sein. Der Mensch sollte also weiterhin in den Prozess eingebunden bleiben, sodass sich Mensch und Maschine gegenseitig verstehen können. Es ist aufzupassen, dass das System nicht zu viele Meldungen rausgibt. Diese sind häufig kaum gefiltert und erhöhen somit den Arbeitsaufwand, der durch die Automatisierung abgenommen werden sollte.
Dies wird auch die \textit{Ironie der Automation} genannt. Während die Automation zu monotonen Aufgaben, wie der Überwachung der Ergebnisse aus den automatisierten Aufgaben, führt, entzieht sie den Nutzer zu stark aus der Aufgabe. Auch hier greift das Konzept des Lernens, indem Aufgaben, die lange Zeit nicht ausgeführt werden, wieder verlernt werden. Sollte das System mit einer Aufgabe mal überfordert sein, kann der Mensch kaum noch weiterhelfen, da dieser den Überblick über das System und dessen Funktion verloren hat und das aufgetretene Problem nicht versteht \cite{Vollrath:2015:Taschenbuch, Wickens:2013}.